% !TeX root = ../../../main.tex
\begin{figure}[htbp]
\begin{ChSixFloatBody}

\centering

\begin{tcolorbox}[ChSix/annopanel]

\begin{minipage}[t]{0.48\textwidth}\vspace{0pt}
 % top-align
\begin{tcolorbox}[
  enhanced,
  colback=blue!3,
  colframe=blue!60!black,
  title=\textbf{PMC3184014},
  fonttitle=\small,
  boxrule=0.8pt,
  arc=2mm,
  left=3mm,right=3mm,top=2mm,bottom=2mm
]
\footnotesize\raggedright
A 39-year-old woman with a history of type 2 diabetes, deep vein thrombosis, gastritis, schizophrenia and self-harm presented to the medical admissions unit with a one-week history of a right sub-mammary abscess followed by spreading cellulitis of the breast. On examination she was septic with a pyrexia of 39.2, blood pressure of 127/65 and tachycardia of 110. Examination of the breast revealed widespread cellulitis involving the nipple and sub-mammary area spreading to the axilla. Initial blood tests showed a white cell count of 23.8 (neutrophils 19.9) and CRP of 428. Simple cellulitis was diagnosed by the admitting physicians and she was commenced on parenteral antibiotics and fluid resuscitation. Her condition worsened despite intravenous benzyl penicillin and flucloxacillin. A surgical opinion was sought on day three when she had evidence of synergistic gangrene. The cellulitis had now spread and there were areas of growing necrotic ulceration (Figure~1). Resuscitation was commenced and antibiotic therapy was adjusted to imipenem and clindamycin based on local necrotizing fasciitis guidelines. She was taken to theatre within a few hours where a partial mastectomy was performed and the wound was left open and packed (Figure~2). Postoperatively she remained stable on intensive care. The following day she returned to theatre for further debridement.

The wound was monitored closely and a vacuum dressing was applied to aid healing. Signs of sepsis improved and she returned to the ward on the fifth postoperative day. Nutritional supplements and adequate hydration were continued during her recovery. Secondary closure was performed on day 13 after the initial operation and there was no further breakdown or compromise of the remaining tissues.

Preoperative cultures taken from the ulcers identified a mixture of gram positive and negative bacteria including \textit{Bacteroides} spp. Histology confirmed widespread microscopic changes and abscess formation consistent with gangrene. There was no evidence of malignancy. She was discharged 22 days after admission and follow-up was arranged.
\end{tcolorbox}
\end{minipage}
\hfill
\begin{minipage}[t]{0.48\textwidth}\vspace{0pt}
 % top-align
\begin{tcolorbox}[ChSix/annobox,title={T2S2 annotation}]
\ChSixAnnotationTableStyle
\begin{tabular}{>{\raggedright\arraybackslash}p{0.78\textwidth} r}
\toprule
\textbf{Event} & \textbf{Time} \\
\midrule
39 years old & 0 \\
woman & 0 \\
type 2 diabetes & -168 \\
deep vein thrombosis & -168 \\
gastritis & -168 \\
schizophrenia & -168 \\
self-harm & -168 \\
presented to the medical admissions unit & 0 \\
right sub-mammary abscess & -168 \\
spreading cellulitis of the breast & -168 \\
septic & 0 \\
pyrexia of 39.2 & 0 \\
blood pressure of 127/65 & 0 \\
tachycardia of 110 & 0 \\
widespread cellulitis involving the nipple & 0 \\
cellulitis of the sub-mammary area & 0 \\
cellulitis spreading to the axilla & 0 \\
white cell count of 23.8 & 0 \\
neutrophils of 19.9 & 0 \\
CRP of 428 & 0 \\
diagnosed with simple cellulitis & 0 \\
commenced on parenteral antibiotics & 0 \\
fluid resuscitation & 0 \\
intravenous benzyl penicillin & 0 \\
flucloxacillin & 0 \\
condition worsened & 72 \\
surgical opinion sought & 72 \\
evidence of synergistic gangrene & 72 \\
cellulitis spread & 72 \\
areas of growing necrotic ulceration & 72 \\
resuscitation commenced & 72 \\
antibiotic therapy adjusted to imipenem & 72 \\
antibiotic therapy adjusted to clindamycin & 72 \\
taken to theatre within a few hours & 72 \\
partial mastectomy performed & 72 \\
wound left open and packed & 72 \\
returned to theatre for further debridement & 96 \\
vacuum dressing applied & 96 \\
signs of sepsis improved & 120 \\
returned to the ward & 120 \\
blood markers improved & 168 \\
secondary closure performed & 312 \\
discharged & 528 \\
follow-up arranged & 528 \\
\bottomrule
\end{tabular}
\end{tcolorbox}
\end{minipage}

\end{tcolorbox}




\end{ChSixFloatBody}
\caption{Clinical narrative (left) aligned with temporal T2S2 annotations (right).}\ChSixFloatLabel{ch06:fig:fine-grained-t2s2-annotation}
\end{figure}